% Insert this content into Section 4.3 (Database Service)

\section{Database Service (db01)}

The database service provides secure data storage with encrypted connections and mutual TLS authentication, implemented as a PostgreSQL 15 server running in a Docker container.

\subsection{Architecture}

\subsubsection{Deployment Model}

The service is deployed as a Docker container (db01) connected only to the internal network (172.19.0.0/16) with static IP 172.19.0.3. This ensures database is not directly accessible from the DMZ, following defense-in-depth principles.

\subsubsection{PostgreSQL Configuration}

\begin{itemize}
    \item \textbf{Version}: PostgreSQL 15
    \item \textbf{Base Image}: postgres:15 official Docker image
    \item \textbf{SSL Mode}: Enabled with client certificate verification
    \item \textbf{Authentication}: SCRAM-SHA-256 + client certificates
    \item \textbf{Port}: 5432 (internal network only)
\end{itemize}

\subsection{Configuration and Deployment}

\subsubsection{TLS Configuration}

PostgreSQL configured to require TLS for all connections:

\textbf{Server Certificates}:
\begin{itemize}
    \item \textbf{Certificate}: db01.org.local-cert.pem (RSA 4096-bit)
    \item \textbf{Private Key}: db01.org.local-key.pem
    \item \textbf{CA Certificate}: ca.crt (Intermediate CA + Root CA chain)
\end{itemize}

\textbf{PostgreSQL SSL Parameters}:
\begin{verbatim}
postgres -c ssl=on \
         -c ssl_cert_file=/var/lib/postgresql/tls-fixed/server.crt \
         -c ssl_key_file=/var/lib/postgresql/tls-fixed/server.key \
         -c ssl_ca_file=/var/lib/postgresql/tls-fixed/ca.crt
\end{verbatim}

\subsubsection{Client Certificate Authentication (mTLS)}

PostgreSQL configured to verify client certificates using pg\_hba.conf:

\begin{verbatim}
# Require SSL with client certificate verification
hostssl all all 0.0.0.0/0 scram-sha-256 clientcert=verify-ca
\end{verbatim}

\textbf{Client Certificate Requirements}:
\begin{itemize}
    \item Certificate must be signed by trusted CA (intermediate-ca)
    \item Must have \texttt{extendedKeyUsage = clientAuth}
    \item Subject CN should match database username (optional)
    \item Presented during TLS handshake
\end{itemize}

\textbf{Client Certificate for web01}:
\begin{itemize}
    \item \textbf{Certificate}: web01-db-client-cert.pem
    \item \textbf{Subject}: CN=web01
    \item \textbf{Key Usage}: Digital Signature, Key Encipherment
    \item \textbf{Extended Key Usage}: clientAuth (critical)
\end{itemize}

\subsubsection{Deployment Steps}

\begin{enumerate}
    \item \textbf{Certificate Generation}: Issue server certificate db01.org.local and client certificate web01-db-client using PKI scripts with proper extendedKeyUsage.
    \item \textbf{Certificate Preparation}: Copy certificates to postgres/tls directory, create entrypoint script to fix permissions (PostgreSQL requires 600 for key files).
    \item \textbf{Database Initialization}: Mount init-users.sql to /docker-entrypoint-initdb.d/ for user creation.
    \item \textbf{Network Configuration}: Assign static IP 172.19.0.3 in internal network only.
    \item \textbf{pg\_hba.conf Configuration}: Configure client certificate verification with clientcert=verify-ca.
    \item \textbf{Container Start}: Launch with \texttt{docker-compose up -d db01}.
    \item \textbf{Configuration Reload}: Apply pg\_hba.conf changes with \texttt{pg\_reload\_conf()}.
\end{enumerate}

\subsection{Security Policies Implementation}

\subsubsection{Encrypted Traffic}

All database connections require TLS:
\begin{itemize}
    \item \textbf{Protocol}: TLS 1.2+ (PostgreSQL default)
    \item \textbf{Cipher Suites}: Strong ciphers (AES-256, ChaCha20)
    \item \textbf{Certificate Verification}: Server presents valid X.509 certificate
    \item \textbf{Client Authentication}: Client presents valid certificate with clientAuth usage
\end{itemize}

\subsubsection{Authentication and Authorization}

\textbf{Multi-Factor Authentication}:
\begin{enumerate}
    \item Password authentication with SCRAM-SHA-256 (hashed passwords)
    \item Client certificate verification (cryptographic proof of identity)
\end{enumerate}

\textbf{Role-Based Access Control (RBAC)}:

\textbf{Status}: Basic roles implemented, extended RBAC planned.

\textbf{Current Implementation}:
\begin{verbatim}
CREATE USER admin WITH PASSWORD 'securepassword';
GRANT ALL PRIVILEGES ON DATABASE securedb TO admin;
\end{verbatim}

\textbf{Planned Extended RBAC}:
\begin{itemize}
    \item \textbf{admin\_user}: Full privileges (CREATE, DROP, INSERT, UPDATE, DELETE, SELECT)
    \item \textbf{app\_user}: Application access (INSERT, UPDATE, SELECT on specific tables)
    \item \textbf{readonly\_user}: Read-only access (SELECT only)
\end{itemize}

Example RBAC configuration:
\begin{verbatim}
-- Create roles
CREATE ROLE readonly_user;
CREATE ROLE app_user;

-- Grant privileges
GRANT SELECT ON ALL TABLES IN SCHEMA public TO readonly_user;
GRANT SELECT, INSERT, UPDATE ON TABLE users TO app_user;

-- Create users with roles
CREATE USER reader WITH PASSWORD 'pass1' IN ROLE readonly_user;
CREATE USER webapp WITH PASSWORD 'pass2' IN ROLE app_user;
\end{verbatim}

\subsubsection{Configuration File Integrity}

\textbf{Critical Configuration Files}:
\begin{itemize}
    \item \texttt{/etc/postgresql/postgresql.conf} (PostgreSQL main configuration)
    \item \texttt{/var/lib/postgresql/data/pg\_hba.conf} (authentication rules)
    \item \texttt{/var/lib/postgresql/tls-fixed/} (TLS certificates and keys)
\end{itemize}

\textbf{Planned Implementation}:
\begin{itemize}
    \item Generate SHA-256 checksums for configuration files
    \item Store checksums in /var/lib/postgresql/config-checksums.sha256
    \item Verify checksums on database startup
    \item Alert on integrity violations
\end{itemize}

\subsubsection{Network Isolation}

Database accessible only from internal network (172.19.0.0/16). Firewall rules allow only web01 (172.19.0.5) → db01 (172.19.0.3):5432. Direct DMZ access blocked by firewall.

\subsection{Database Schema}

\textbf{Users Table}:
\begin{verbatim}
CREATE TABLE users (
    id SERIAL PRIMARY KEY,
    username VARCHAR(255) UNIQUE NOT NULL,
    password_hash VARCHAR(255) NOT NULL,
    totp_secret VARCHAR(255),
    created_at TIMESTAMP DEFAULT CURRENT_TIMESTAMP
);
\end{verbatim}

\subsection{Implementation Decisions}

\subsubsection{mTLS Client Certificates}

\textbf{Decision}: Require client certificates with clientAuth extended key usage.

\textbf{Justification}: Provides cryptographic proof of identity beyond passwords. PostgreSQL strictly enforces certificate purpose - serverAuth certificates rejected for client authentication. Discovered during implementation: initial client certificate had serverAuth, causing "certificate verify failed" error.

\subsubsection{Internal Network Only}

\textbf{Decision}: Database connected only to internal network, not DMZ.

\textbf{Justification}: Database contains sensitive data and should not be directly accessible from public-facing network. All access mediated through web01 application layer. Firewall explicitly blocks DMZ → database connections.

\subsubsection{SCRAM-SHA-256 Authentication}

\textbf{Decision}: Use SCRAM-SHA-256 for password hashing instead of MD5.

\textbf{Justification}: SCRAM-SHA-256 provides strong password hashing with salt and iteration count, protecting against rainbow table attacks. MD5 is cryptographically broken and should not be used for security-critical functions.

\subsection{Client Certificate Generation Script}

Created dedicated script \texttt{PKI/scripts/issue-client-cert.sh} for generating client certificates with proper clientAuth usage:

\begin{verbatim}
#!/bin/bash
# Generate client certificate with clientAuth extended key usage
cat > "$TEMP_EXT_CONFIG" << EOF
basicConstraints = CA:FALSE
keyUsage = critical, digitalSignature, keyEncipherment
extendedKeyUsage = clientAuth
subjectKeyIdentifier = hash
authorityKeyIdentifier = keyid,issuer
EOF

openssl ca -config openssl.cnf -days 365 -md sha256 \
  -in client.csr -out client-cert.pem \
  -extfile "$TEMP_EXT_CONFIG"
\end{verbatim}

This script prevents the serverAuth/clientAuth confusion encountered during initial implementation.
